\section{Chapter~\ref{sec:debugging}. Debugging the machine}

What an electrode hears, the low dimensionality of activity, the performance of interfaces on rigid arrays, the joint learning of brain and decoder, and the visual spots produced by stimulation are measured in peer-reviewed experiments; the data-rate estimates are the chapter's arithmetic; for the Neuralink device implanted in humans, as far as we could check, data have been published mostly by the company itself.

{\footnotesize
\setlength{\tabcolsep}{3pt}\renewcommand{\arraystretch}{1.15}
\begin{longtable}{>{\raggedright}p{0.5cm}>{\raggedright}p{4.0cm}>{\raggedright}p{5.6cm}>{\raggedright}p{2.3cm}>{\raggedright\arraybackslash}p{1.7cm}}
\toprule
No. & Claim & Experiment and what was measured & Source & Status \\
\midrule\endfirsthead
\toprule
No. & Claim & Experiment and what was measured & Source & Status \\
\midrule\endhead
1 & An electrode hears the spikes of neurons within a radius of about a hundred micrometers, usually one to several; they are told apart by shape & Rat, area CA1, simultaneous intracellular and extracellular recording: the extracellular spike shape follows the width and amplitude of the intracellular one. Estimate: a tetrode could separate $60$--$100$ neurons, but in practice about six are isolated. & \cite{Henze2000extra} & measured \\
2 & The brain has $8.6\cdot10^{10}$ neurons; a probe hears about one hundred-millionth & Counting nuclei in a suspension of dissolved tissue: $86\cdot10^9$ neurons in the human brain. The fraction is the chapter's arithmetic. & \cite{Azevedo2009} & measured \\
3 & The activity of hundreds of motor-cortex neurons fits into one or two dozen shared variables & Review of recordings in monkey motor cortex: population activity is described by a few latent variables shared by many neurons. & \cite{Gallego2017} & measured \\
4 & Neighbors' noise is correlated, and a hundred neurons carry almost as much as a thousand (Proposition~\ref{prop:pooling}) & Monkey visual cortex: the noise correlation coefficient of neighboring neurons is about $0.1$, and the gain from averaging saturates at a hundred neurons. Theory of information-limiting correlations. & \cite{Zohary1994, MorenoBote2014} & measured \\
5 & Raw stream of $2\cdot10^8$~bit/s versus $8\cdot10^4$~bit/s in spikes~\eqref{eq:probedata} & The chapter's arithmetic from the capacity of the spike stream~\eqref{eq:spikeentropy}. The digitization parameters are real: the Neuralink chip samples at $19.3$~kHz, $10$~bits per channel. & derived in chapter; \cite{Rieke1997, Musk2019} & theory \\
6 & The first Neuralink system: chips amplify and digitize, the raw stream goes over a cable, spikes are detected by software outside; on-chip detection, a sealed housing, radio, and wireless power belong to the device for humans, according to the company & Company paper: the full raw stream over a USB-C cable, spikes detected by real-time software off the chip; on-board compression and wireless power and transmission are named as a plan for clinical devices. For the device implanted in humans, only company reports. That on-chip spike detection is necessary is the chapter's conclusion from~\eqref{eq:probedata}. & \cite{Musk2019}; Neuralink reports, no peer-reviewed paper & measured (company paper); for humans, company report \\
7 & Up to $3072$ electrodes on $96$ threads of $32$; a robot inserts the threads, avoiding vessels & Company paper: $3072$ electrodes on $96$ threads of $32$; the robot inserts $6$ threads ($192$ electrodes) per minute; up to $70\,\%$ of electrodes in rats with a chronically implanted array record spikes. & \cite{Musk2019} & measured (company paper) \\
8 & In humans, about a thousand channels on $64$ threads; some threads shifted; cursor at about $10$~bit/s & Company reports only. A PubMed search found no peer-reviewed paper with human results; this agrees with the caveat in the chapter text. & Neuralink reports; no peer-reviewed paper & measured (company report) \\
9 & An interface on an array of a hundred electrodes controls a cursor & A person with paralysis, $96$ electrodes in primary motor cortex: the intention to move the arm modulates spikes three years after the injury; the decoder gives a ``neural cursor'' (email, television), control of a prosthetic hand and a robotic arm. & \cite{Hochberg2006} & measured \\
10 & Writing with the ``mental hand'': about $90$ characters per minute, less than $8$~bit/s & A person with a paralyzed hand, attempted handwriting, recurrent decoder: $90$ characters per minute at $94.1\,\%$ accuracy in real time, above $99\,\%$ after autocorrection. The estimate in bits is the chapter's arithmetic. & \cite{Willett2021} & measured \\
11 & Attempts to speak are turned into text at about $60$ words per minute & A person who lost intelligible speech, arrays in the cortex: $62$ words per minute; $9.1\,\%$ word errors with a $50$-word vocabulary, $23.8\,\%$ with a $125\,000$-word vocabulary. & \cite{Willett2023} & measured \\
12 & An interface extracts a small fraction of the capacity: tens of bits per second for one neuron, thousands for a hundred (Proposition~\ref{prop:probe}) & The upper bound~\eqref{eq:spikeentropy} is theory; the comparison with rows~8 and 10 is the chapter's conclusion. That the bottleneck is low dimensionality and ignorance of the code rather than the wire has not been tested directly. & \cite{Rieke1997} & theory \\
13 & The brain and the decoder learn from each other & Monkeys control a cursor; the decoder adapts in closed loop while the neurons relearn at the same time: accuracy rises, the skill is retained and suffers less from movements of the monkey's own arm. The advice to separate the learning rates is an engineering rule; the chapter has no separate experiment for it. & \cite{Orsborn2014coad} & measured \\
14 & Current through an electrode excites the surrounding neurons and wires regardless of type & Mouse and cat cortex, two-photon calcium imaging: even a weak current excites sparse neurons scattered over up to millimeters, through direct excitation of axons in a volume tens of micrometers across. So the excitation is nonselective but not solid. & \cite{Histed2009stim} & measured \\
15 & Stimulation of the visual cortex lights spots; up to $16$ spots at once allow recognizing some letters and object edges & A blind person, $96$ electrodes in the visual cortex for $6$ months: single-electrode threshold $66.8\pm36.5$~\textmu A; simultaneous stimulation of groups (up to $16$ electrodes, the stimulator's limit) lowers the threshold and gives distinguishable patterns; some letters and object edges are recognized. & \cite{Fernandez2021} & measured \\
16 & Neuralink's second direction is a device that feeds signals into the visual cortex & Only company reports of development; no peer-reviewed data on its performance were found. & Neuralink reports & hypothesis \\
17 & Concentrations of broadcast transmitters can be measured with a chemical sensor in the tissue & Rat brain slices, fast-scan cyclic voltammetry at a carbon fiber: serotonin release and reuptake measured; the substance spills beyond the synapse. & \cite{Bunin1998} & measured \\
\bottomrule
\end{longtable}}
